% !TEX root = ../main.tex
%-----------------------------------------------------------------------
%  Title block and abstract.
%  \input by ../main.tex -- not compiled on its own.
%-----------------------------------------------------------------------
\title{\bfseries Parametric Adjoint Sensitivities in the DINO
       Configuration of MITgcm\\[0.45em]
       \large Part~I: a vertical-diffusivity perturbation ensemble\\
       Part~II: a neural-network surrogate for the sensitivity fields}
\author{Tanvir Shahriar}
% No reader is named here. These documents are shared, and a name in the
% title block dates the copy and presumes on the reader.
\date{5 August 2026 \quad$\cdot$\quad revised 21 August 2026}
\maketitle

\begin{abstract}
\noindent
The MITgcm--Tapenade adjoint of the idealised DINO basin
($51\times198\times36$) returns spatially resolved sensitivities of a meridional
transport cost functional to the vertical diffusivity field $\kv(x,y,z)$, at
about $6\times10^{2}$ core-hours for a five-year sweep. Those sensitivities hold
at one set of parameter values, and whether they remain valid when the
parameters change is not known for this configuration.

This document sets out a two-part programme. \textbf{Part~I} is a
seven-member perturbation ensemble varying $\kv$ over a factor of $128$, which
asks directly whether the sensitivity patterns depend on the mixing
parameterisation. It is costed from measured run times, carries a resource
request, and is now running. \textbf{Part~II} works through what it would take
to replace the adjoint calculation with a trained network mapping
\emph{(background ocean state, parameter field, lead time, cost-functional
specification)} to the sensitivity field, so that parametric sensitivity costs
an evaluation rather than an integration.

The two are sequential. Part~I is a scoping study whose primary output is a
design answer --- whether mixing must be an input to the surrogate at all ---
and whose training pairs are a by-product. Its result determines what Part~II
should be. Its second result, a finite-difference check taken where the
sensitivity is large, would be the first meaningful validation this adjoint has
had.
\end{abstract}
